\begin{abstract}
Grounded drug question answering requires both retrieving the appropriate evidence and avoiding unsupported answers, yet Vietnamese lacks dedicated resources for evaluating these capabilities. We construct a Vietnamese drug-information corpus from the national drug formulary and package leaflets, together with a 10{,}000-question retrieval benchmark whose answer locations are known by construction and a 500-item end-to-end evaluation set covering single-turn, multi-turn, unanswerable, out-of-scope, and prompt-injection requests. On this benchmark, hybrid retrieval with a self-hosted reranker places a correct passage in the top ten for 97.5\% of questions. On top of this retriever, we assemble an agentic retrieval-augmented generation pipeline for low-cost language models from established components and limit every model call to one small structured decision. With the commercial Gemini 3.8 Flash it reaches a truthfulness of 0.923 and reduces hallucinated answers eighteenfold, and the open-weight Qwen3.5-9B within the same pipeline closely matches it. Beyond providing the first dedicated benchmark for Vietnamese drug-information retrieval, our results reveal that grounding in the corpus accounts for most of the accuracy gain and that each agentic step pays off on the questions it targets, with query rephrasing answering conversational follow-ups that one-step retrieval leaves unanswered and outperforming the simple concatenation of earlier questions in both model settings.
\keywords{\sloppy\hyphenpenalty=10000 \exhyphenpenalty=10000 Retrieval-augmented generation \and Agentic RAG \and Drug information retrieval \and Pharmaceutical documents \and Retrieval benchmark \and LLM-as-a-judge}
\end{abstract}
